\section{Limitations and Future Work}
\label{sec:limitations}
Our evaluation is currently restricted to multiple-choice benchmarks, and the binary flip metric may not transfer to open-ended or free-form settings where sycophancy manifests differently. Future work should extend the framework to open-ended generation tasks, where sycophancy can be measured through semantic drift rather than answer flips. HLE coverage is currently partial: pressure pipelines have been completed for GPT-5.4 and Gemini 3.5 Flash, but the remaining four HLE cells (Claude Haiku, Claude Sonnet, GPT-5.4 Mini, GPT-5.4 Nano) are deferred to future work (Table~\ref{tab:coverage}). Cross-model HLE claims, in particular the fragmenter mode, are therefore based on a two-model contrast and should be read as a stress-test signal rather than a fully cross-family generalization. The reasoning-trace analysis relies on an external model to calibrate belief confidence across turns, which introduces the calibrator's confound of own errors and biases entangled with the signal being measured, and its accuracy on adversarial traces is not independently validated. A self-consistency calibration approach, or validation against human annotations on a held-out trace set, would decouple the calibration signal from the target model's behavior. The corrigibility gap metric distinguishing sycophantic capitulation from legitimate belief updating is defined but not fully operationalized. Building a parallel set of genuine evidence-based counter-arguments would enable direct measurement of how well models discriminate valid correction from social pressure.
